% List of symbols, read by thesis.tex in the preamble.
%   \defsymbol[unit]{key}{template}[name in the list]{description}
% Write \sym{key} in the text or in a formula to print the symbol with a link to its row of the list; only the symbols
% used at least once appear in the list, in the order in which they are defined here.
% Subscripts and superscripts go right after the key and stay inside the link: \sym{sigma}_k^2, \sym{bezier}_i^{\prime}.
% In the template, \symsub is "_{index}" when a subscript is given (\mu\symsub: \sym{mu}_i is \mu_i), while \symsep and
% \symidx add the index to a subscript of the symbol itself (\lambda_{s\symsep\symidx}: \sym{lams}_i is \lambda_{s,i}).
% Changing a template changes the symbol everywhere; every key is one meaning. The only letters left with two meanings
% are p (distribution and position), q (density q_theta and waypoints q_j), f (state update f_theta and waypoint frame)
% and C (continuity class C^0 and Cholesky factors), which the context and the subscripts tell apart.
% Chapter 4 does not use this list yet.

% ------------------------------------------------------------------ kinematic model
\defsymbol[s]{t}{t\symsub}{time; $t_n$ is the timestamp of sample $n$}
\defsymbol[s]{tzero}{t_{0\symsub}}{onset time of a stroke}
\defsymbol{mu}{\mu\symsub}{log-time delay of a stroke}
\defsymbol{sigma}{\sigma\symsub}{log-response time of a stroke}
\defsymbol[m]{amplitude}{D\symsub}{amplitude of a stroke, the length of the path it describes}
\defsymbol[s$^{-1}$]{lognormal}{\Lambda\symsub}{lognormal speed profile of unit area}
\defsymbol[s$^2$]{tvar}{V\symsub}{variance of the lognormal speed profile of a stroke, seen as a distribution over time}
\defsymbol{logtime}{\chi\symsub}{logarithm of the time elapsed since the onset of a stroke, $\ln(t-t_0)$}
\defsymbol[m/s]{velocity}{\vec v\symsub}{velocity of the end effector; $\vec v_n$ is the measured velocity of sample $n$}
\defsymbol{nsegments}{M}{number of strokes of a movement; in a recording, the number of kept strokes (segments)}

% ------------------------------------------------------------------ Bézier model
\defsymbol[m]{bezier}{B\symsub}{cubic Bézier curve; $B_i$ is the curve of segment $i$}
\defsymbol[m]{ctrlpt}{P\symsub}[P_0,\dots,P_3]{control points of a Bézier curve: $P_0$ and $P_3$ are its end points, $P_1$ and $P_2$ its inner control points; $P_{j,i}$ belongs to segment $i$}
\defsymbol{curvepar}{s}{parameter of a Bézier curve, in $[0,1]$}
\defsymbol{czero}{C^0}{continuity of a spline: consecutive curves share their end points}
\defsymbol{gone}{G^1}{geometric continuity of a spline: consecutive curves also share their tangent direction}
\defsymbol{gratio}{\kappa\symsub}{ratio between the tangent vectors of segments $i$ and $i+1$ at their junction}
\defsymbol{direction}{u\symsub}{unit tangent direction at a node of the chain}
\defsymbol{direstimate}{\hat u\symsub}{finite-difference estimate of a tangent direction}
\defsymbol[rad]{rotation}{\alpha\symsub}{rotation of a tangent direction from its estimate}
\defsymbol{rotbasis}{o\symsub}[(o_{j,1},o_{j,2})]{orthonormal basis of the plane orthogonal to $\hat u_j$}
\defsymbol[m]{lams}{\lambda_{s\symsep\symidx}}{start tangent length of a segment}
\defsymbol[m]{lame}{\lambda_{e\symsep\symidx}}{end tangent length of a segment}
\defsymbol[m]{arclength}{L\symsub}{arc length of a segment}
\defsymbol{normarc}{\ell\symsub}{normalised arc length along a Bézier curve}
\defsymbol{dataarc}{\xi\symsub}{normalised arc length of a recorded sample within its segment}

% ------------------------------------------------------------------ data and extraction
\defsymbol{nsamples}{N}{number of samples of a recording}
\defsymbol[m]{position}{p\symsub}{measured position; $p_n$ is the position of sample $n$}
\defsymbol[s]{samplingint}{\Delta t}{sampling interval}
\defsymbol{minimum}{n\symsub}{sample index of a speed minimum}
\defsymbol{nwindows}{M_0}{number of windows between the speed minima of a recording, before the strokes are filtered}
\defsymbol{window}{\mathcal W\symsub}{window of the samples of a stroke}
\defsymbol[s]{dtb}{\Delta t_b}{minimum time between two speed minima}
\defsymbol[m/s]{vb}{v_b}{minimum prominence of a speed minimum}
\defsymbol[m/s]{residual}{r\symsub}[r^{(k)}]{residual speed before the fit of stroke $k$}
\defsymbol{mumin}{\mu_{\min}}{lower bound of $\mu$ in the lognormal fit}
\defsymbol{mumax}{\mu_{\max}}{upper bound of $\mu$ in the lognormal fit}
\defsymbol{sigmamin}{\sigma_{\min}}{lower bound of $\sigma$ in the lognormal fit}
\defsymbol{sigmamax}{\sigma_{\max}}{upper bound of $\sigma$ in the lognormal fit}
\defsymbol[m/s]{vpeak}{v^{\mathrm{peak}}\symsub}{peak speed of a stroke}
\defsymbol[m/s]{vmin}{v_{\min}}{peak-speed threshold of the kept strokes}
\defsymbol[m]{displacement}{\Delta p\symsub}{displacement of a stroke, from its start point to its end point}
\defsymbol[m]{lammin}{\lambda_{\min}}{lower bound of the tangent lengths}
\defsymbol[m]{lammax}{\lambda_{\max\symsep\symidx}}[\lambda_{\max,i}]{upper bound of the tangent lengths of segment $i$}
\defsymbol[rad]{alphamax}{\alpha_{\max}}{largest rotation of a tangent direction}
\defsymbol[s]{tpeakwin}{\hat t}{time of the largest residual speed in the window of a stroke}
\defsymbol{onsetfactor}{\gamma}{onset factor of a starting point of the lognormal fit}
\defsymbol{ssigma}{\mathcal S_\sigma}{starting values of $\sigma$ of the lognormal fit}
\defsymbol{sonset}{\mathcal S_\gamma}{onset factors of the starting points of the lognormal fit}
\defsymbol{ngl}{n_{\mathrm{GL}}}{nodes of the Gauss--Legendre quadrature}
\defsymbol{nell}{n_\ell}{grid values of the curve parameter used to invert $\ell_i$}
\defsymbol{epsl}{\varepsilon_L}{relative tolerance on the length constraint}

% ------------------------------------------------------------------ generative models
\defsymbol{obs}{x\symsub}{observation; $x_t$ is element $t$ of a sequence, and $x_i$ the output features of stroke $i$}
\defsymbol{obsspace}{\mathcal X}{space of the observations}
\defsymbol{dimension}{d}{dimension of an observation; for the network, the number of output dimensions}
\defsymbol{condition}{y}{condition of a generative model}
\defsymbol{condspace}{\mathcal Y}{space of the conditions}
\defsymbol{pdata}{p_{\mathrm{data}}}{distribution of the data}
\defsymbol{pdist}{p}{probability distribution, or its density}
\defsymbol{pmodel}{p_\theta}{distribution defined by a model with parameters $\theta$}
\defsymbol{family}{\mathcal P_\Theta}{parametric family of distributions}
\defsymbol{decisionfn}{F_\theta}{decision function of a discriminative model}
\defsymbol{modelparams}{\theta}{parameters of a model}
\defsymbol{paramspace}{\Theta}{space of the parameters of a model}
\defsymbol{divergence}{\mathcal D}{divergence between two distributions}
\defsymbol{seqlen}{T}{length of a sequence}
\defsymbol{condstep}{\phi\symsub}{part of the condition relevant to step $t$; $\phi_i$ are the conditioning inputs of stroke $i$}
\defsymbol{markovorder}{J}{order of a Markov chain}
\defsymbol{state}{h\symsub}{state of a recurrent model; $h_i$ is the hidden state of the LSTM at stroke $i$}
\defsymbol{statefn}{f_\theta}{state update of a recurrent model}
\defsymbol{qtheta}{q_\theta}{conditional density of one step of an autoregressive model, here a Gaussian mixture}
\defsymbol{ncomp}{K}{number of components of a Gaussian mixture}
\defsymbol{mixweight}{\pi\symsub}{weight of a mixture component}
\defsymbol{mixmean}{m\symsub}{mean of a mixture component}
\defsymbol{mixchol}{C\symsub}{lower-triangular (Cholesky) factor of the covariance of a mixture component}
\defsymbol{covariance}{\Sigma\symsub}{covariance of a mixture component}
\defsymbol{gaussian}{\mathcal N}{Gaussian (normal) density}

% ------------------------------------------------------------------ features of the network
\defsymbol[m]{waypoint}{q\symsub}{waypoint of a task}
\defsymbol[m]{wpvec}{w\symsub}{vector from the start of a stroke to the waypoint it aims at}
\defsymbol[m]{wpnext}{w'\symsub}{vector from the start of a stroke to the following waypoint}
\defsymbol{refnormal}{\nu\symsub}{unit normal of the reference plane of a stroke}
\defsymbol{wpframe}{f\symsub}[(f_1,f_2,f_3)]{waypoint frame, $f_1$ along $w_i$}
\defsymbol{strokeframe}{e\symsub}[(e_1,e_2,e_3)]{stroke frame, $e_1$ along the chord $\Delta p_i$}
\defsymbol{chorddir}{\hat d\symsub}{unit direction of the chord of a stroke}
\defsymbol[s]{tpeak}{t^{\mathrm{peak}}\symsub}{time of the speed peak of a stroke}
\defsymbol{peakdelay}{\tau\symsub}{log of the delay between the speed peaks of two consecutive strokes, relative to $e^{\mu_{i-1}}$}
\defsymbol{extent}{\rho\symsub}{log of the length of the chord relative to the waypoint vector}
\defsymbol[rad]{azimuth}{\psi\symsub}{direction of the chord within the reference plane}
\defsymbol[rad]{elevation}{\beta\symsub}{direction of the chord out of the reference plane}
\defsymbol{handle}{\eta\symsub}{tangent handle in the stroke frame, relative to the chord and compressed by $\operatorname{asinh}$}
\defsymbol{handlescale}{\lambda_\eta}{scale of the inverse hyperbolic sine of the handles, relative to the chord}
\defsymbol[s]{tfirst}{t_0^{\mathrm{first}}}{onset of the first stroke of a generated sequence}

% ------------------------------------------------------------------ network and training
\defsymbol{previn}{\tilde x\symsub}{previous-stroke inputs: the features of stroke $i-1$ given to the LSTM at stroke $i$}
\defsymbol{lstmin}{z\symsub}{input of the LSTM at a stroke}
\defsymbol{cellstate}{c\symsub}{cell state of the LSTM}
\defsymbol{trunkout}{g\symsub}{output of the trunk layer}
\defsymbol{headhidden}{a\symsub}{output of the hidden layer of the head}
\defsymbol{weights}{W\symsub}{weight matrix of a layer}
\defsymbol{bias}{b\symsub}{bias vector of a layer}
\defsymbol{rawfactor}{\tilde C\symsub}{free entries of a Cholesky factor, as computed by the head}
\defsymbol{diagfloor}{\delta}{lower bound of the diagonal of the Cholesky factors}
\defsymbol{determinism}{\omega}{determinism of the sampling, in $[0,1]$}
\defsymbol{stdnoise}{\zeta}{standard normal vector used to sample a mixture component}
\defsymbol{loss}{\mathcal L}{training loss, the mean negative log-likelihood per stroke}
\defsymbol{nbatch}{N_b}{number of strokes of a batch}
\defsymbol{dimset}{S\symsub}{output dimensions that exist for a stroke}

% ------------------------------------------------------------------ evaluation
\defsymbol{predsample}{\hat x\symsub}{output features of a stroke sampled from the predicted density $q_\theta$}
\defsymbol[m]{tubepoint}{\tilde p\symsub}{point of a path aligned by normalized arc length; $\tilde p_{k,m}$ is point $m$ of segment $k$}
\defsymbol[m]{tubemean}{\bar p\symsub}{mean of the aligned recorded paths at a point of the tube}
\defsymbol[m$^2$]{tubecov}{\Gamma\symsub}{covariance of the aligned recorded paths at a point of the tube}
\defsymbol{tubebasis}{\Omega\symsub}{orthonormal basis of the plane orthogonal to the mean path at a point of the tube}
\defsymbol{mahal}{d_{\mathrm M\symsep\symidx}}{Mahalanobis distance of a point of a path from the recorded paths}
\defsymbol{tailprob}{Q\symsub}{tail probability of the Mahalanobis distance at a point of a path}
\defsymbol{pathscore}{Q_{\min}}{score of a path, the smallest tail probability over its points}
\defsymbol{tubethreshold}{Q^\ast}{threshold of the whole-path tube}
